\ncertchapter{3}{Vectors}\label{ch:vectors}

One dimension is enough to define velocity and acceleration, but it is not enough to describe almost
any real motion. A thrown ball moves sideways and up at once; a car turns a corner; a planet sweeps
round the Sun. To handle these we need quantities that carry a direction as well as a size, and rules
for combining them that respect the direction. Those quantities are vectors, and this chapter is about
the arithmetic they obey.

The chapter's central claim is that \emph{the single most useful skill in mechanics is resolving a
vector into perpendicular components}. Resolution turns one
hard two-dimensional problem into two easy one-dimensional ones, each of which can be handled with the
tools of \chref{ch:straightline}. Almost every later chapter relies on it, from projectiles to inclined planes to
banked roads, and a student who resolves fluently will find that most of mechanics is easier than it
looks.

Before resolution, though, comes a more basic question that is often skipped: what actually makes a
quantity a vector? Having a direction is not enough, and the chapter begins by saying why.

\section{What Makes a Quantity a Vector}

A \textbf{scalar} (mass, time, temperature, energy) is fully specified by a number and a unit. A
\textbf{vector} (displacement, velocity, force) also needs a \emph{direction}, and, crucially, it
must add by the triangle rule (\S\ref{sec:vecadd}): two walks, \SI{3}{m} east then \SI{4}{m} north, take
you \SI{5}{m} from the start, not \SI{7}{m}. In print a vector is written in
bold with an arrow, $\vc{A}$; its magnitude, a positive scalar with the same units, is written $|\vc{A}|$
or simply $A$.

\subsection{The vocabulary: types of vectors}

Vector problems are more often lost in the wording than the algebra, so it is worth fixing the names
once. Each of the following is standard exam vocabulary; Figure~\ref{fig:vectypes} draws four of them
side by side.

\begin{defbox}{Types of vectors}
\begin{itemize}
\item \textbf{Unit vector} $\uhat{a}=\vc{a}/|\vc{a}|$ (read ``$a$-hat''): magnitude exactly $1$ and no
      units, carrying direction only. The unit vectors along the $x$, $y$ and $z$ axes are written
      $\ihat,\jhat,\khat$.
\item \textbf{Null (zero) vector} $\vc{0}$: zero magnitude, indeterminate direction; it is what you get from
      $\vc{a}-\vc{a}$, or from the displacement of a body that returns to its start.
\item \textbf{Equal vectors:} same magnitude \emph{and} direction, regardless of where they are drawn.
\item \textbf{Negative vector} $-\vc{a}$: same magnitude, exactly reversed direction.
\item \textbf{Parallel vectors:} along the same or parallel lines (\emph{collinear} when drawn along
      one line), \emph{like} if they
      point the same way, \emph{unlike} if opposite.
\item \textbf{Coplanar vectors:} all lying in one plane. Any two vectors are coplanar; three need
      not be.
\item \textbf{Co-initial vectors:} drawn from a common starting point.
\item \textbf{Position vector} $\vr$: from a chosen origin to a point, locating it.
      \textbf{Displacement vector:} from an initial to a final position, $\Delta\vr=\vr_2-\vr_1$.
\item \textbf{Axial vectors:} represent rotation and point along the rotation axis by the right-hand
      rule (\S\ref{sec:vecprod}; examples are angular velocity $\vomega$, torque $\vtau$ and angular momentum $\vL$, all met in
      \chref{ch:rotation}), as opposed to ordinary \emph{polar} vectors like force and
      velocity.
\end{itemize}
\end{defbox}

\begin{figure}[htbp]
\centering
\tikzfig{vector_types}
\figcaption{fig:vectypes}{Four of the types drawn. Equal vectors need not start from the same point;
a negative vector is a reversal, not a subtraction; collinear vectors share a line; coplanar vectors
share a plane.}
\end{figure}

\begin{trapbox}{a direction alone does not make a vector}
Electric \emph{current} has a direction but is a \textbf{scalar}: currents meeting at a junction add
arithmetically, not by the triangle law. The test for a vector is not ``does it have a direction'' but
``does it add like a displacement''. Finite rotations fail the same test: a rotation of $90^\circ$
about $x$ followed by $90^\circ$ about $y$ leaves a book in a different orientation from the reverse
order, so finite rotations do not add commutatively and are not vectors. Infinitesimal rotations do
commute, which is why angular velocity $\vomega$ is a vector.
\end{trapbox}

\section{Adding Vectors, and Resolving Them}\label{sec:vecadd}

The defining property of a vector is the way it adds. There are two standard constructions, and they always
agree.

\begin{keybox}{Two equivalent additions}
\textbf{Tip-to-tail (triangle):} place $\vc{B}$'s tail at $\vc{A}$'s tip; the resultant runs from $\vc{A}$'s tail to
$\vc{B}$'s tip. \textbf{Parallelogram:} draw $\vc{A}$ and $\vc{B}$ from a common origin; the resultant is the
diagonal. They give the same vector, as drawing both on the same pair (Figure~\ref{fig:vectoradd}) shows.
\end{keybox}

\begin{figure}[htbp]
\centering
\tikzfig{kin_vectoradd}
\figcaption{fig:vectoradd}{The triangle and parallelogram constructions give identical resultants. For magnitudes, the law of cosines gives $R=\sqrt{A^2+B^2+2AB\cos\theta}$, where $\theta$ is the angle between the vectors drawn tail to tail; in the triangle it is the exterior angle between $\vc{A}$ produced and $\vc{B}$, not the triangle's interior angle.}
\end{figure}

\begin{figure}[htbp]
\centering
\tikzfig{vec_parallelogram}
\figcaption{fig:vecpara}{The construction behind the resultant formula. Dropping the perpendicular $CN$
onto the extended baseline gives a right triangle whose legs are $B\sin\theta$ (rise) and
$A+B\cos\theta$ (run); Pythagoras on that triangle is the whole derivation, and the same triangle gives
the direction $\phi$.}
\end{figure}

\begin{derivbox}{The magnitude of the resultant}
Let $\theta$ be the angle between $\vc{A}$ and $\vc{B}$ drawn tail to tail. Place $\vc{A}$ along a
baseline and draw $\vc{B}$ from the head of $\vc{A}$, at angle $\theta$ to the extended baseline
(Figure~\ref{fig:vecpara}). Drop a perpendicular from the head of $\vc{B}$
(the point $C$) to the extended baseline, meeting it at $N$. Then the horizontal run is $A+B\cos\theta$ and the vertical
rise is $B\sin\theta$. By Pythagoras the resultant length is
\[
R=\sqrt{(A+B\cos\theta)^2+(B\sin\theta)^2}=\sqrt{A^2+B^2+2AB\cos\theta},
\]
using $\cos^2\theta+\sin^2\theta=1$. Its direction $\phi$ above $\vc{A}$ satisfies
$\tan\phi=\dfrac{B\sin\theta}{A+B\cos\theta}$. Two checks: $\theta=0$ gives $R=A+B$ (parallel vectors
just add); $\theta=180^\circ$ gives $R=|A-B|$ (opposed vectors subtract).
\end{derivbox}

\subsection{Subtraction, and adding many vectors}

Subtraction is addition of a reversed vector, and adding any number of vectors is repeated tip-to-tail
addition. Both follow at once from the triangle rule.

\begin{keybox}{Vector subtraction and the polygon law}
\textbf{Subtraction:} $\vc{A}-\vc{B}=\vc{A}+(-\vc{B})$: reverse $\vc{B}$ and add. Drawn from a common tail, $\vc{A}-\vc{B}$ runs from
the \emph{head of $\vc{B}$ to the head of $\vc{A}$}. \\[3pt]
\textbf{Scalar multiple:} $k\vc{A}$ has magnitude $|k|A$ and points along $\vc{A}$ if $k>0$, against it
if $k<0$. In particular every vector is its magnitude times its unit vector, $\vc{A}=A\,\uhat{A}$. \\[3pt]
\textbf{Order does not matter:} $\vc{A}+\vc{B}=\vc{B}+\vc{A}$ (the two routes round a parallelogram), and
grouping does not matter either, $(\vc{A}+\vc{B})+\vc{C}=\vc{A}+(\vc{B}+\vc{C})$. \\[3pt]
\textbf{Polygon law:} to add several vectors, place them head-to-tail in turn; the resultant closes the
polygon, from the first tail to the last head. If the vectors themselves close the polygon (last head
meets first tail), their sum is the null vector.
\end{keybox}

\begin{defbox}{The equilibrant}
The \textbf{equilibrant} of a set of vectors is the single vector that restores equilibrium: equal in
magnitude to their resultant $\vc{R}$, but \emph{opposite} in direction. Adding it makes the total zero.
This is the entire basis of force equilibrium: a particle is in equilibrium precisely when the forces on
it form a closed polygon, so their resultant, and the force needed to balance them, is zero. We
use exactly this in \chref{ch:newton} (Lami's theorem and the triangle of forces).
\end{defbox}

\subsection{Components: the master technique}

Geometric addition is clear on paper and fragile in practice. Resolving every vector into perpendicular
components replaces the geometry with arithmetic, and it works for any number of vectors.

\begin{defbox}{Resolution into components}
Any vector $\vc{A}$ in a plane equals the sum of its projections on two perpendicular axes:
\[
\vc{A} = A_x\,\ihat + A_y\,\jhat,\qquad A_x=A\cos\theta,\quad A_y=A\sin\theta,
\]
where $\theta$ is the angle from the positive $x$-axis to $\vc{A}$.
The magnitude is $A=\sqrt{A_x^2+A_y^2}$ and direction $\tan\theta=A_y/A_x$; the signs of $A_x$ and $A_y$
decide the quadrant, since $\tan\theta$ alone cannot tell $(1,1)$ from $(-1,-1)$. In three dimensions a
third component is added: $\vc{A}=A_x\ihat+A_y\jhat+A_z\khat$ and $A=\sqrt{A_x^2+A_y^2+A_z^2}$. The
components are scalars (with sign), the unit vectors carry the directions. If $\alpha,\beta,\gamma$ are
the angles $\vc{A}$ makes with the three axes, then $A_x=A\cos\alpha$, $A_y=A\cos\beta$, $A_z=A\cos\gamma$,
and these \textbf{direction cosines} obey $\cos^2\alpha+\cos^2\beta+\cos^2\gamma=1$ (divide
$A^2=A_x^2+A_y^2+A_z^2$ by $A^2$).
\end{defbox}

\begin{methodbox}{Adding vectors by components}
\stepn{1} Resolve every vector along the same pair of perpendicular axes, taking care over signs.
\stepn{2} Add the components axis by axis: $R_x=A_x+B_x+C_x+\cdots$ and $R_y=A_y+B_y+C_y+\cdots$. This
is legitimate because $\ihat$ and $\jhat$ are fixed, so the $\ihat$ terms and the $\jhat$ terms can be
collected separately.
\stepn{3} Recombine: $R=\sqrt{R_x^2+R_y^2}$, with direction $\tan\phi=R_y/R_x$ and the quadrant read
from the signs of $R_x$ and $R_y$.
\end{methodbox}

\begin{exbox}{Three displacements}
A walker goes \SI{4.0}{km} east, then \SI{6.0}{km} in a direction $60^\circ$ north of east, then
\SI{3.0}{km} west. Find the resultant displacement.

Take $x$ east and $y$ north. The components are $(4.0,\,0)$, $(6.0\cos60^\circ,\,6.0\sin60^\circ)=(3.0,\,5.2)$
and $(-3.0,\,0)$, in kilometres. Adding, $R_x=4.0+3.0-3.0=\SI{4.0}{km}$ and $R_y=\SI{5.2}{km}$, so
\[
R=\sqrt{4.0^2+5.2^2}=\SI{6.6}{km},\qquad \phi=\tan^{-1}\frac{5.2}{4.0}=52^\circ\ \text{north of east}.
\]
The triangle construction would need the resultant formula twice; components needed only addition.
\ans{\SI{6.6}{km} at $52^\circ$ north of east}
\end{exbox}

\inwords{Resolving a vector splits one slanted quantity into a horizontal part and a vertical part that
can be handled separately. Whenever the acceleration along $x$ does not depend on the motion along $y$
(as for a projectile without air resistance), one two-dimensional problem becomes two one-dimensional
problems, and one-dimensional problems we can already solve.}

\begin{trapbox}{$\cos$ goes with the angle you measure from}
$A\cos\theta$ is the component \emph{along} the axis from which $\theta$ is measured, and $A\sin\theta$
is the component perpendicular to it. Swapping them is the single most common vector error. Always ask:
``which axis is my angle measured from?''
\end{trapbox}

\section{Multiplying Vectors}\label{sec:vecprod}

Addition combines like with like, but some quantities depend on two directions at once: the work done
by a force depends on how much of it lies \emph{along} the motion, the turning effect of a force on how
much lies \emph{across} the line to the pivot. The two products previewed in \S\ref{sec:vecalg}
answer these questions: the \textbf{dot product} $\vc{A}\cdot\vc{B}=AB\cos\theta$, a scalar, and the
\textbf{cross product} $\vc{A}\ex\vc{B}$, a vector of magnitude $AB\sin\theta$, where $\theta$
($0\le\theta\le180^\circ$) is the angle between them drawn tail to tail. Now that vectors can be
written in components, we can find each product's component form and see why it is right.

\begin{keybox}{Properties of the dot product}
$B\cos\theta$ is the component of $\vc{B}$ along $\vc{A}$, so $\vc{A}\cdot\vc{B}=A\times$(component of
$\vc{B}$ along $\vc{A}$), and the component of any vector along a unit vector $\uhat{n}$ is
$\vc{A}\cdot\uhat{n}$. The product is commutative, $\vc{A}\cdot\vc{B}=\vc{B}\cdot\vc{A}$, and distributive,
$\vc{A}\cdot(\vc{B}+\vc{C})=\vc{A}\cdot\vc{B}+\vc{A}\cdot\vc{C}$. It is positive for an acute angle, zero
for a right angle, negative for an obtuse one, and $\vc{A}\cdot\vc{A}=A^2$. For the axes,
$\ihat\cdot\ihat=\jhat\cdot\jhat=\khat\cdot\khat=1$ and $\ihat\cdot\jhat=\jhat\cdot\khat=\khat\cdot\ihat=0$.
\end{keybox}

\begin{derivbox}{Why $\vc{A}\cdot\vc{B}=A_xB_x+A_yB_y+A_zB_z$}
Write both vectors in components and multiply out using the distributive law. Of the nine terms, the six
that pair two \emph{different} unit vectors vanish, since the axes are perpendicular; the three that pair
a unit vector with itself carry a factor $1$. What survives is $A_xB_x+A_yB_y+A_zB_z$. Equating this to
$AB\cos\theta$ gives the working formula for the angle between two vectors,
$\cos\theta=(A_xB_x+A_yB_y+A_zB_z)/AB$.
\end{derivbox}

\begin{keybox}{Direction of the cross product}
$\vc{A}\ex\vc{B}$ is perpendicular to the plane of $\vc{A}$ and $\vc{B}$. Its sense is given by the
\textbf{right-hand rule}: point the fingers of the right hand along $\vc{A}$ and curl them towards
$\vc{B}$ through the smaller angle; the thumb points along $\vc{A}\ex\vc{B}$. Reversing the order
reverses the curl, so $\vc{B}\ex\vc{A}=-\vc{A}\ex\vc{B}$, and $\vc{A}\ex\vc{A}=\vc{0}$ (so
$\ihat\ex\ihat=\vc{0}$, and likewise for $\jhat$, $\khat$). For the axes,
$\ihat\ex\jhat=\khat$, $\jhat\ex\khat=\ihat$, $\khat\ex\ihat=\jhat$ (cyclic order gives $+$, the
reverse order $-$). Multiplying out the nine terms as for the dot product gives
\[
\vc{A}\ex\vc{B}=(A_yB_z-A_zB_y)\,\ihat+(A_zB_x-A_xB_z)\,\jhat+(A_xB_y-A_yB_x)\,\khat
=\begin{vmatrix}\ihat&\jhat&\khat\\ A_x&A_y&A_z\\ B_x&B_y&B_z\end{vmatrix},
\]
the determinant being a compact way to remember it (expand along the top row; the $\jhat$ term takes a
minus sign). The rule is drawn in Figure~\ref{fig:cross}, where
torque first needs it.
\end{keybox}

\begin{exbox}{Angle, projection and area for one pair}
Take $\vc{A}=3\ihat+4\jhat$ and $\vc{B}=4\ihat+3\jhat$, so $A=B=5$. The dot product is
$\vc{A}\cdot\vc{B}=12+12=24$, so $\cos\theta=24/25$ and $\theta=16.3^\circ$. The unit vector along
$\vc{B}$ is $\uhat{B}=0.8\ihat+0.6\jhat$, and the component of $\vc{A}$ along $\vc{B}$ is
$\vc{A}\cdot\uhat{B}=2.4+2.4=4.8$. The cross product is $\vc{A}\ex\vc{B}=(3\cdot3-4\cdot4)\khat=-7\khat$:
the parallelogram they span has area $7$ (check: $AB\sin\theta=25\times0.28=7$), and the minus sign
says turning from $\vc{A}$ to $\vc{B}$ is clockwise seen from $+z$. \ans{$16.3^\circ$; $4.8$; $-7\khat$}
\end{exbox}

\begin{trapbox}{there is no dividing by a vector}
$\vc{A}\cdot\vc{B}=0$ with neither vector zero means they are perpendicular, not that one vanishes.
Likewise $\vc{A}\cdot\vc{B}=\vc{A}\cdot\vc{C}$ does \emph{not} give
$\vc{B}=\vc{C}$, only that $\vc{B}-\vc{C}$ is perpendicular to $\vc{A}$.
\end{trapbox}

The dot product returns in \chref{ch:energy} as work, $W=\vF\cdot\vc{d}$, and the cross product in
\chref{ch:rotation} as torque and angular momentum; the problem types below practise both.

\section{Summary}\label{sec:summary3}
A vector is a quantity with magnitude and direction that \emph{also} adds by the triangle law; a
direction alone does not qualify, which is why electric current, despite having a sense, is a scalar.
Vectors add tip to tail, or equivalently by the parallelogram rule, and subtraction is addition of the
reversed vector. Several vectors add by the polygon law, and if they close the polygon their sum is zero,
which is precisely the condition for a particle to be in equilibrium.

\begin{itemize}
\item Resultant of two vectors at angle $\theta$: $R=\sqrt{A^2+B^2+2AB\cos\theta}$, always between
      $|A-B|$ and $A+B$.
\item Resolution: $A_x=A\cos\theta$, $A_y=A\sin\theta$, with $\theta$ measured \emph{from the $x$-axis};
      on a slope inclined at $\theta$, the weight's component along the incline is $mg\sin\theta$, because
      $\theta$ is the angle between the weight and the normal.
\item For three or more vectors, resolve, add components, and recombine. This never needs the resultant
      formula and works for any number of vectors.
\item The equilibrant is equal and opposite to the resultant; three forces in equilibrium form a closed
      triangle.
\item $\vc{A}\cdot\vc{B}=AB\cos\theta=A_xB_x+A_yB_y+A_zB_z$, a scalar, gives the angle between two
      vectors and the component of one along the other.
\item $\vc{A}\ex\vc{B}$ has magnitude $AB\sin\theta$ (the parallelogram area), direction by the
      right-hand rule, and components $(A_yB_z-A_zB_y,\ A_zB_x-A_xB_z,\ A_xB_y-A_yB_x)$, from
      $\ihat\ex\jhat=\khat$ and its cyclic partners.
\end{itemize}

\section{Problem Types}\label{sec:types3}
Vector questions come in five recurring shapes, and every one of them is settled most reliably by
resolving into components rather than by geometry: components are mechanical, and mechanical is what
you want when time is short.

\archq{Resultant of two vectors, or the angle between them}
{``Two forces/velocities of given magnitude act at a given angle, find the resultant,'' or the reverse:
``their resultant has magnitude $R$, find the angle.''}
{Use $R=\sqrt{A^2+B^2+2AB\cos\theta}$ and $\tan\phi=\dfrac{B\sin\theta}{A+B\cos\theta}$. For
perpendicular vectors this reduces to Pythagoras. Remember $R$ always lies between $|A-B|$ and $A+B$, which
alone kills most wrong options.}
{Adding magnitudes arithmetically ($6+8=14$) instead of vectorially. Also: $|\vc A+\vc B|=|\vc A-\vc B|$
does \emph{not} mean the vectors are equal; for non-zero vectors it means they are perpendicular.}
{Two forces of \SI{6}{N} and \SI{8}{N} act at right angles to each other. The magnitude of their
resultant is:}
{\SI{2}{N}}{\SI{10}{N}}{\SI{14}{N}}{\SI{48}{N}\archans{B}{Perpendicular, so
$R=\sqrt{6^2+8^2}=\sqrt{100}=\SI{10}{N}$. Option (C) is the trap of adding magnitudes; (A) is their
difference.}}

\archq{Angle between two vectors from the dot product}
{``Find the angle between $\vc A$ and $\vc B$,'' or ``show they are perpendicular.''}
{Use $\cos\theta=\dfrac{\vc A\cdot\vc B}{|\vc A||\vc B|}$, with $\vc A\cdot\vc B$ computed from
components. A zero dot product (with neither vector zero) means \textbf{perpendicular}; a zero cross
product means \textbf{parallel} or antiparallel.}
{Forgetting to divide by the magnitudes, or assuming a negative dot product is an error; it
means only that the angle is obtuse.}
{The angle between $\vc A=\ihat+\jhat$ and $\vc B=\ihat-\jhat$ is:}
{$0^\circ$}{$45^\circ$}{$90^\circ$}{$180^\circ$\archans{C}{$\vc A\cdot\vc B=(1)(1)+(1)(-1)=0$, and
neither vector is zero, so they are perpendicular.}}

\archq{Maximum, minimum and possible resultants}
{``Two forces of given magnitudes act at an unknown angle, which resultant is possible/impossible?''}
{The resultant is bounded: $|A-B|\le R\le A+B$. Maximum when parallel ($\theta=0$), minimum when
antiparallel ($\theta=180^\circ$). Anything outside that band is impossible, whatever the angle.}
{Assuming any value is achievable by choosing the angle. The band is fixed by the two magnitudes alone.}
{Two forces of \SI{5}{N} and \SI{3}{N} act at a point. Which resultant is \emph{impossible}?}
{\SI{2}{N}}{\SI{5}{N}}{\SI{8}{N}}{\SI{9}{N}\archans{D}{The resultant must lie between $|5-3|=2$ and
$5+3=8$. \SI{9}{N} lies outside, so it cannot occur for any angle.}}

\archq{Cross product, area and the determinant}
{``Find $\vc A\ex\vc B$,'' or ``find the area of the parallelogram/triangle formed by two vectors.''}
{Evaluate the determinant
$\vc A\ex\vc B=\begin{vmatrix}\ihat&\jhat&\khat\\ A_x&A_y&A_z\\ B_x&B_y&B_z\end{vmatrix}$.
Its magnitude is the \textbf{parallelogram} area; for a \emph{triangle} take half. Always check the
result is perpendicular to both by dotting it with each.}
{Sign errors in the $\jhat$ term; it carries a minus in the cofactor expansion. And the cross product
is \emph{anti}-commutative: $\vc A\ex\vc B=-\vc B\ex\vc A$.}
{The area of the triangle formed by $\vc A=2\ihat$ and $\vc B=3\jhat$ is:}
{$3$}{$6$}{$12$}{$\sqrt{13}$\archans{A}{$\vc A\ex\vc B=6\khat$, so the parallelogram area is $6$ and the
triangle is half of that: $3$ square units.}}

\archq{Resolving a vector, and the equilibrium of three forces}
{``Find the components of a vector,'' or ``three forces keep a body in equilibrium, find one of
them.''}
{Resolve everything into $x$ and $y$: $A_x=A\cos\theta$, $A_y=A\sin\theta$ (checking which is adjacent to
the marked angle). For equilibrium set $\sum F_x=0$ \emph{and} $\sum F_y=0$; equivalently the three
vectors form a closed triangle, so Lami's theorem applies.}
{Pairing $\cos$ with the horizontal automatically. That holds only when $\theta$ is measured from the
$x$-axis; on an incline the along-slope component is $mg\sin\theta$.}
{Three forces in equilibrium: $\vc F_1=3\ihat+4\jhat$ and $\vc F_2=-\ihat+2\jhat$ (in newtons). The
third force $\vc F_3$ is:}
{$2\ihat+6\jhat$}{$-2\ihat-6\jhat$}{$4\ihat+2\jhat$}{$-4\ihat-2\jhat$\archans{B}{Equilibrium means
$\vc F_1+\vc F_2+\vc F_3=\vc0$, so $\vc F_3=-(\vc F_1+\vc F_2)=-(2\ihat+6\jhat)=-2\ihat-6\jhat$.}}
